\section{Conclusion}
\label{sec:conclusion}

This article formalizes Euclid's theorem from the same first-principles
foundation used throughout the preceding chapters. The proof follows the
classical primorial-plus-one construction: from a finite list of primes it
builds a number that is congruent to one modulo every prime in the list,
then uses the existence of a smallest divisor to extract a prime factor
outside that list. The contradiction is mathematical before it is
computational: no member of the original list can divide the constructed
number, while the constructed number must still have a prime divisor.

The Stainless development verifies each step that the article relies on:
small-remainder facts, zero-remainder preservation under multiplication,
divisibility of product members, smallest-divisor primality, and the
final non-membership theorem. The result is a source-backed formal proof
of the infinitude of primes, with the supporting finite-prefix corollaries
separated from the theorem spine rather than folded into the main claim.

The theorem spine can be read compactly as the following verified chain.
Let $P=\operatorname{primorial}(L)$ and $N=P+1$:

\begin{equation*}
\begin{aligned}
\forall p\in L,\quad \operatorname{mod}(N,p)&=1\ne0
&&\text{[Stage 1]} \\
d=\operatorname{findSmallestDivisor}(N,2),\quad d=N
&\Rightarrow \operatorname{isPrime}(N) \\
d<N
&\Rightarrow \operatorname{isPrime}(d)
&&\text{[Stage 2]} \\
v\mid N\ \land\ p\in L
&\Rightarrow p\ne v
&&\text{[Stage 3].}
\end{aligned}
\end{equation*}

Thus either $N$, or its least non-trivial divisor $d$, is a prime outside
$L$:

\begin{equation*}
L\ne[]\ \Rightarrow\ \exists p:\operatorname{isPrime}(p)\land p\notin L.
\quad\blacksquare\ \text{[Q.E.D.]}
\end{equation*}

\section{Future Work}
\label{sec:future-work}

The most natural continuation is the Fundamental Theorem of Arithmetic,
since Euclid's theorem already establishes the existence side of prime
decomposition. A verified uniqueness proof would require a stronger
library of divisibility and coprimality lemmas, but it would extend the
present result in a direct and structurally compatible way.

Further work could then move from existence to distribution. Dirichlet's
theorem would require arithmetic progressions and substantially richer
modular reasoning, while the Prime Number Theorem would require asymptotic
analysis far beyond the finite arithmetic developed here. Those
directions are intentionally outside the scope of this article, but this
proof supplies a verified starting point for them.
